% GENERATED FILE. Edit canonical lessons, then run npm run generate:books.
% canonical-source-hash: fnv1a64:4c20511ad667dd37
% canonical-lessons: TA-C557-parutti, TA-C557-kampali, TA-C557-cayam, TA-C557-kittar, TA-C557-vayalin

\chapter{Cloth and Instruments for a Fair}
\label{ch:ta-cloth-and-instruments-for-a-fair}
\hlchaptermodality{557}

\begin{chapteropening}
\textbf{By the end of this chapter:} \emph{I can talk about cloth, wool and Western instruments.}

Complete the last lesson of chapter 557: I can talk about cloth, wool and Western instruments.
\end{chapteropening}

\section[\texorpdfstring{\ta{பருத்தி} (parutti)}{பருத்தி (parutti)}]{\ta{பருத்தி} (parutti) — cotton}

\label{lesson:TA-C557-parutti}



\begin{quote}
Before the new one: say the Tamil for weaving, then the Tamil for gold thread.
\end{quote}

\subsection*{You'll want to know: \ta{பருத்தி}}
\textbf{\ta{பருத்தி}} — \emph{parutti} — \textquotedblleft{}cotton\textquotedblright{}.

\begin{quote}
\ta{கோடையில்} \ta{பருத்தித்} \ta{துணி} \ta{அணிந்தால்} \ta{வியர்க்காது}.
\end{quote}

\begin{quote}
\emph{kōṭaiyil paruttit tuṇi aṇintāl viyarkkātu.}
\end{quote}

\begin{quote}
If you wear cotton clothes in summer, you will not sweat.
\end{quote}

Before a noun it doubles the next stop: \ta{பருத்திச்} \ta{சேலை}, \ta{பருத்தித்} \ta{துணி}.

\begin{cousinweb}
\textbf{\ta{பருத்தி}} (\emph{parutti}) is native Dravidian: DEDR 3976 gives \ta{பருத்தி} (\emph{parutti}), \textquotedblleft{}the Indian cotton plant; cotton\textquotedblright{}.
\end{cousinweb}

\begin{grammarlens}[title={What you've built}]
One more word for saying what happens if.
\end{grammarlens}

\subsection*{Your turn}
\begin{itemize}
\raggedright
  \item \emph{Say it:} \emph{parutti}
  \item \emph{Say it:} \emph{parutti}, once more
  \item \emph{Say it:} \emph{kōṭaiyil paruttit tuṇi aṇintāl viyarkkātu.} — If you wear cotton clothes in summer, you will not sweat.
  \item \emph{Recall:} say the Tamil for weaving, then the Tamil for gold thread, then say \emph{parutti} again
\end{itemize}

\begin{tcolorbox}[breakable,colback=teal!4,colframe=teal!35!black,title={Before you move on}]
What does \ta{பருத்தி} mean? (\textquotedblleft{}Cotton\textquotedblright{}.) Say it once more.
\end{tcolorbox}
\section[\texorpdfstring{\ta{கம்பளி} (kampaḷi)}{கம்பளி (kampaḷi)}]{\ta{கம்பளி} (kampaḷi) — wool; a woollen blanket}

\label{lesson:TA-C557-kampali}



\begin{quote}
Before the new one: say the Tamil for gold thread, then the Tamil for cotton.
\end{quote}

\subsection*{You'll want to know: \ta{கம்பளி}}
\textbf{\ta{கம்பளி}} — \emph{kampaḷi} — \textquotedblleft{}wool; a woollen blanket\textquotedblright{}.

\begin{quote}
\ta{குளிர்} \ta{அதிகமானால்} \ta{தாத்தா} \ta{கம்பளி} \ta{போர்த்திக்கொள்வார்}.
\end{quote}

\begin{quote}
\emph{kuḷir atikamāṉāl tāttā kampaḷi pōrttikkoḷvār.}
\end{quote}

\begin{quote}
If it gets very cold, Grandfather wraps himself in a woollen blanket.
\end{quote}

Woollen clothes are \ta{கம்பளி} \ta{ஆடைகள்}.

\begin{cousinweb}
\textbf{\ta{கம்பளி}} (\emph{kampaḷi}) is borrowed from Sanskrit \emph{kambala}, \textquotedblleft{}a woollen blanket or upper garment\textquotedblright{} (CDIAL 2771).
\end{cousinweb}

\begin{grammarlens}[title={What you've built}]
One more word for saying what happens if.
\end{grammarlens}

\subsection*{Your turn}
\begin{itemize}
\raggedright
  \item \emph{Say it:} \emph{kampaḷi}
  \item \emph{Say it:} \emph{kampaḷi}, once more
  \item \emph{Say it:} \emph{kuḷir atikamāṉāl tāttā kampaḷi pōrttikkoḷvār.} — If it gets very cold, Grandfather wraps himself in a woollen blanket.
  \item \emph{Recall:} say the Tamil for gold thread, then the Tamil for cotton, then say \emph{kampaḷi} again
\end{itemize}

\begin{tcolorbox}[breakable,colback=teal!4,colframe=teal!35!black,title={Before you move on}]
What does \ta{கம்பளி} mean? (\textquotedblleft{}Wool; a woollen blanket\textquotedblright{}.) Say it once more.
\end{tcolorbox}
\section[\texorpdfstring{\ta{சாயம்} (cāyam)}{சாயம் (cāyam)}]{\ta{சாயம்} (cāyam) — a dye; colour}

\label{lesson:TA-C557-cayam}



\begin{quote}
Before the new one: say the Tamil for cotton, then the Tamil for wool.
\end{quote}

\subsection*{You'll want to know: \ta{சாயம்}}
\textbf{\ta{சாயம்}} — \emph{cāyam} — \textquotedblleft{}a dye; colour\textquotedblright{}.

\begin{quote}
\ta{மழையில்} \ta{நனைந்தால்} \ta{இந்தத்} \ta{துணியின்} \ta{சாயம்} \ta{போகும்}.
\end{quote}

\begin{quote}
\emph{maḻaiyil naṉaintāl intat tuṇiyiṉ cāyam pōkum.}
\end{quote}

\begin{quote}
If it gets wet in the rain, this cloth will lose its dye.
\end{quote}

Colour coming out of cloth is \ta{சாயம்} \ta{போ}.

\begin{cousinweb}
\textbf{\ta{சாயம்}} (\emph{cāyam}) is native Dravidian: DEDR 2457 gives \ta{சாய்} (\emph{cāy}), \textquotedblleft{}brilliance, colour\textquotedblright{}, and \ta{சாயம்} (\emph{cāyam}), \textquotedblleft{}colour, tint, dye\textquotedblright{}.
\end{cousinweb}

\begin{grammarlens}[title={What you've built}]
One more word for saying what happens if.
\end{grammarlens}

\subsection*{Your turn}
\begin{itemize}
\raggedright
  \item \emph{Say it:} \emph{cāyam}
  \item \emph{Say it:} \emph{cāyam}, once more
  \item \emph{Say it:} \emph{maḻaiyil naṉaintāl intat tuṇiyiṉ cāyam pōkum.} — If it gets wet in the rain, this cloth will lose its dye.
  \item \emph{Recall:} say the Tamil for cotton, then the Tamil for wool, then say \emph{cāyam} again
\end{itemize}

\begin{tcolorbox}[breakable,colback=teal!4,colframe=teal!35!black,title={Before you move on}]
What does \ta{சாயம்} mean? (\textquotedblleft{}A dye; colour\textquotedblright{}.) Say it once more.
\end{tcolorbox}
\section[\texorpdfstring{\ta{கிட்டார்} (kiṭṭār)}{கிட்டார் (kiṭṭār)}]{\ta{கிட்டார்} (kiṭṭār) — a guitar}

\label{lesson:TA-C557-kittar}



\begin{quote}
Before the new one: say the Tamil for wool, then the Tamil for a dye.
\end{quote}

\subsection*{You'll want to know: \ta{கிட்டார்}}
\textbf{\ta{கிட்டார்}} — \emph{kiṭṭār} — \textquotedblleft{}a guitar\textquotedblright{}.

\begin{quote}
\ta{நீ} \ta{கிட்டார்} \ta{கற்றுக்கொண்டால்} \ta{நாம்} \ta{சேர்ந்து} \ta{பாடலாம்}.
\end{quote}

\begin{quote}
\emph{nī kiṭṭār kaṟṟukkoṇṭāl nām cērntu pāṭalām.}
\end{quote}

\begin{quote}
If you learn the guitar, we can sing together.
\end{quote}

Played with \ta{வாசி}: \ta{கிட்டார்} \ta{வாசி}.

\begin{cousinweb}
\textbf{\ta{கிட்டார்}} (\emph{kiṭṭār}) is a loan from English \emph{guitar}.
\end{cousinweb}

\begin{grammarlens}[title={What you've built}]
One more word for saying what happens if.
\end{grammarlens}

\subsection*{Your turn}
\begin{itemize}
\raggedright
  \item \emph{Say it:} \emph{kiṭṭār}
  \item \emph{Say it:} \emph{kiṭṭār}, once more
  \item \emph{Say it:} \emph{nī kiṭṭār kaṟṟukkoṇṭāl nām cērntu pāṭalām.} — If you learn the guitar, we can sing together.
  \item \emph{Recall:} say the Tamil for wool, then the Tamil for a dye, then say \emph{kiṭṭār} again
\end{itemize}

\begin{tcolorbox}[breakable,colback=teal!4,colframe=teal!35!black,title={Before you move on}]
What does \ta{கிட்டார்} mean? (\textquotedblleft{}A guitar\textquotedblright{}.) Say it once more.
\end{tcolorbox}
\section[\texorpdfstring{\ta{வயலின்} (vayaliṉ)}{வயலின் (vayaliṉ)}]{\ta{வயலின்} (vayaliṉ) — a violin}

\label{lesson:TA-C557-vayalin}



\begin{quote}
Before the new one: say the Tamil for a dye, then the Tamil for a guitar.
\end{quote}

\subsection*{You'll want to know: \ta{வயலின்}}
\textbf{\ta{வயலின்}} — \emph{vayaliṉ} — \textquotedblleft{}a violin\textquotedblright{}.

\begin{quote}
\ta{அவர்} \ta{வயலின்} \ta{வாசித்தால்} \ta{எல்லோரும்} \ta{அமைதியாகக்} \ta{கேட்பார்கள்}.
\end{quote}

\begin{quote}
\emph{avar vayaliṉ vācittāl ellōrum amaitiyākak kēṭpārkaḷ.}
\end{quote}

\begin{quote}
When he plays the violin, everyone listens quietly.
\end{quote}

Not a form of \ta{வயல்} (\emph{vayal}), a field: a loan from English.

\begin{cousinweb}
\textbf{\ta{வயலின்}} (\emph{vayaliṉ}) is a loan from English \emph{violin}.
\end{cousinweb}

\begin{grammarlens}[title={What you've built}]
One more word for saying what happens if.
\end{grammarlens}

\subsection*{Your turn}
\begin{itemize}
\raggedright
  \item \emph{Say it:} \emph{vayaliṉ}
  \item \emph{Say it:} \emph{vayaliṉ}, once more
  \item \emph{Say it:} \emph{avar vayaliṉ vācittāl ellōrum amaitiyākak kēṭpārkaḷ.} — When he plays the violin, everyone listens quietly.
  \item \emph{Recall:} say the Tamil for a dye, then the Tamil for a guitar, then say \emph{vayaliṉ} again
\end{itemize}

\begin{tcolorbox}[breakable,colback=teal!4,colframe=teal!35!black,title={Before you move on}]
What does \ta{வயலின்} mean? (\textquotedblleft{}A violin\textquotedblright{}.) Say it once more.
\end{tcolorbox}
